\section{Conclusiones}

\begin{itemize}
    \item La implementación de \textbf{Odoo 19} ha permitido a \textbf{Machu Picchu Exclusive Tours E.I.R.L.} consolidar la información de sus clientes y reservas en una plataforma única, reduciendo la dispersión de datos en canales informales.
    \item Se ha fortalecido el \textbf{control interno} sobre los adelantos y saldos, disminuyendo el riesgo asociado a la compra de servicios no reembolsables sin una validación previa adecuada.
    \item El uso de tableros Kanban y actividades ha contribuido a \textbf{ordenar el flujo operativo} de cada viaje, facilitando la priorización de tareas y el seguimiento de hitos clave.
    \item La organización de la información y la mejora en tiempos de respuesta sientan las bases para un \textbf{crecimiento escalable} de la agencia en el segmento de turismo receptivo de lujo.
    \item El ecosistema configurado se encuentra alineado con una estrategia de \textbf{digitalización progresiva}, permitiendo en el futuro la incorporación de nuevos módulos (contabilidad, facturación electrónica, e--commerce) sin perder coherencia con lo ya implementado.
\end{itemize}

